\documentclass[a4paper,11pt]{article}

\usepackage[utf8]{inputenc}
\usepackage[T1]{fontenc}
\usepackage{amsmath,amssymb,bm}
\usepackage{geometry}
\usepackage{hyperref}
\geometry{margin=2.2cm}
\hypersetup{
  colorlinks=true,
  linkcolor=blue,
  urlcolor=blue
}

\title{A Brief Methodology Overview for High-Precision Data Processing\\(Extrapolation, Uncertainty Propagation, Fitting, and Statistics)}
\author{}
\date{}

\begin{document}
\maketitle

\section{Notation and Basic Assumptions}

This document only provides a high-level theoretical overview of the methods. It does not describe any specific UI or software implementation details.

\subsection{Data and uncertainty}

Let the measured quantity be $x$ and its standard uncertainty be $\sigma_x$ (standard uncertainty). In uncertainty propagation and weighted fitting, we assume different inputs are independent (uncorrelated) by default. If correlations exist, a covariance-matrix formulation should be used (see below).

\subsection{Sequences and limits}

The typical object of extrapolation is a convergent sequence $\{s_n\}$ (e.g., an energy sequence as the basis size / truncation level increases). The goal is to estimate the limit $s_\infty=\lim_{n\to\infty}s_n$ and provide a credible uncertainty estimate.

\section{Sequence Extrapolation Methods (Extrapolation)}

\subsection{Three-point extrapolation (default three-point formula)}

Given three values $A,B,C$ (often viewed as three progressively improved approximations), a commonly used three-point extrapolation is
\begin{equation}
V = \frac{(C-B)^2}{B-A} + C,
\end{equation}
where $V$ is the extrapolated value. Key properties:
\begin{itemize}
  \item Only three points are required; no extra model parameters.
  \item Invariant under additive shifts of the data, but numerically unstable when $B\approx A$ (denominator near zero).
  \item Useful as a fast and conservative three-point limit estimate, but not guaranteed to be valid for all sequence types.
\end{itemize}

A conservative uncertainty estimate can be taken as the maximum deviation from the input points:
\begin{equation}
U=\max\left(|V-A|,|V-B|,|V-C|\right).
\end{equation}
This $U$ is closer to a consistency-difference measure than a statistical standard deviation.

\subsection{Power-law extrapolation}

For many numerical approximations (especially sequences associated with truncation order / basis size), a common leading-error model is
\begin{equation}
E(x) = E_\infty + a\,x^{-p},
\end{equation}
where $x$ controls accuracy (e.g., basis level or truncation index), $E_\infty$ is the target limit, and $a,p$ are unknown constants. Given three points $(x_1,E_1),(x_2,E_2),(x_3,E_3)$, one can eliminate variables to solve for $p$ and obtain $E_\infty$:
\begin{align}
R &= \frac{E_1-E_2}{E_2-E_3},\\
\frac{x_1^{-p}-x_2^{-p}}{x_2^{-p}-x_3^{-p}} &= R,\\
a &= \frac{E_1-E_2}{x_1^{-p}-x_2^{-p}},\qquad
E_\infty = E_1 - a\,x_1^{-p}.
\end{align}
Notes:
\begin{itemize}
  \item The model has a clear numerical/physical motivation (dominant error decays as a power law).
  \item Three-point data can be sensitive: if $E_2\approx E_3$ or $x_2^{-p}\approx x_3^{-p}$, the solution may become unstable.
  \item Robustness can be improved if $p$ is known or constrained; multi-point fitting generalizes this to an overdetermined problem.
\end{itemize}

A common conservative uncertainty estimate is the difference to a reference value (often the highest-level point $E_3$ or a user-selected reference column):
\begin{equation}
U = |E_\infty - E_{\mathrm{ref}}|.
\end{equation}

\subsection{Richardson extrapolation / sequence acceleration}

The core idea of Richardson extrapolation: if an approximation $S(h)$ has an asymptotic expansion
\begin{equation}
S(h)=S + c\,h^p + \mathcal{O}(h^{p+1}),
\end{equation}
then combining results at different step sizes $h_1,h_2$ can eliminate the leading error term and increase the convergence order. When $p$ is known, for example,
\begin{equation}
S \approx \frac{h_2^p S(h_1) - h_1^p S(h_2)}{h_2^p - h_1^p}.
\end{equation}
With more points, an extrapolation table can be constructed to recursively eliminate higher-order terms. Notes:
\begin{itemize}
  \item Works well when the leading error follows a smooth asymptotic form with a clear order $p$.
  \item Can be unstable for sequences without a smooth asymptotic expansion, or with oscillatory / non-analytic errors.
  \item Requires some knowledge of the effective order $p$ (or an adaptive estimate).
\end{itemize}

In practice, when a reliable statistical error model is absent, a conservative estimate often uses a reference term $S_{\mathrm{ref}}$ (e.g., the last term or a specified approximation) and sets
\begin{equation}
U=|S-S_{\mathrm{ref}}|.
\end{equation}
This reflects the difference between the accelerated value and the reference term, and does not necessarily equal a statistical standard deviation; cross-checks via differences and multiple methods are recommended.

\subsection{Shanks transform and Wynn $\varepsilon$ algorithm}

The Shanks transform and Wynn $\varepsilon$ algorithm are nonlinear sequence transformations, often used to accelerate linearly convergent or logarithmically convergent sequences. Intuitively, they approximate the limit via rational approximants (closely related to Pad\'e ideas).

The simplest accelerated form is related to Aitken's $\Delta^2$ process:
\begin{equation}
s'_n = s_n - \frac{(\Delta s_n)^2}{\Delta^2 s_n},\quad
\Delta s_n=s_{n+1}-s_n,\quad
\Delta^2 s_n=s_{n+2}-2s_{n+1}+s_n,
\end{equation}
which can significantly accelerate linear convergence when $\Delta^2 s_n$ is not near zero. Wynn's $\varepsilon$ algorithm provides a systematic recursive table construction; in practice one typically uses the recursion directly rather than writing formulas by hand.

Notes:
\begin{itemize}
  \item Often effective for linearly convergent sequences and closely connected to Pad\'e approximants.
  \item Can suffer from division-by-zero or numerical amplification when differences are small or the sequence is noisy.
  \item For strongly oscillatory sequences it may be less robust than Levin-type methods.
\end{itemize}

\subsection{Levin transforms (with $u$-transform as a representative)}

Levin-type transforms introduce a remainder estimate $\omega_n$ to construct a better estimate of the limit, and are commonly used for slowly convergent or alternating/oscillatory sequences. A general form is
\begin{equation}
S \approx \frac{\sum_{k=0}^{n} w_k\,\frac{s_k}{\omega_k}}{\sum_{k=0}^{n} w_k\,\frac{1}{\omega_k}},
\end{equation}
where the weights $w_k$ depend on the chosen variant ($u,t,v$, etc.), and $\omega_k$ is usually derived from differences (e.g., $s_{k+1}-s_k$) or a heuristic model.

Notes:
\begin{itemize}
  \item Often outperforms Shanks/Wynn for alternating or slowly convergent sequences.
  \item Depends on the quality of $\omega_k$: when the remainder estimate degenerates (e.g., identical consecutive terms leading to ``zero weights''), additional handling is required.
  \item Still sensitive to noise; diagnostic quantities and cross-method comparisons are recommended.
\end{itemize}

\section{Uncertainty Propagation Methods (Uncertainty Propagation)}

\subsection{First-order (linearized) propagation}

Let the output be
\begin{equation}
y=f(x_1,\ldots,x_m),
\end{equation}
When input uncertainties are small and $f$ is smooth near the operating point, a first-order Taylor expansion yields an approximate standard uncertainty:
\begin{equation}
\sigma_y^2 \approx \sum_{i=1}^{m}\left(\frac{\partial f}{\partial x_i}\sigma_{x_i}\right)^2.
\end{equation}
If inputs are correlated, a covariance-matrix form is used:
\begin{equation}
\sigma_y^2 \approx \bm{g}^\mathsf{T}\,\bm{\Sigma}\,\bm{g},\qquad
\bm{g}=\nabla f,\ \bm{\Sigma}_{ij}=\mathrm{Cov}(x_i,x_j).
\end{equation}
Typical notes:
\begin{itemize}
  \item Suitable as a local approximation for ``small uncertainty + smooth function''.
  \item May underestimate uncertainty for strong nonlinearity, non-Gaussian inputs, or large uncertainties; higher-order or Monte Carlo methods are recommended.
\end{itemize}

\subsection{Obtaining partial derivatives: analytic/symbolic vs numerical differences}

First-order propagation requires partial derivatives $\partial f/\partial x_i$. In principle:
\begin{itemize}
  \item Analytic derivatives: manual derivation or symbolic differentiation to obtain exact expressions.
  \item Numerical derivatives: when analytic forms are hard, approximate with finite differences.
\end{itemize}

Central differences provide a common second-order approximation:
\begin{equation}
\frac{\partial f}{\partial x}\bigg|_{x=x_0}\approx \frac{f(x_0+h)-f(x_0-h)}{2h}.
\end{equation}
The step size $h$ must balance truncation error (decreases as $h$ decreases) and rounding error (increases as $h$ decreases). For second-order central differences, an empirical choice is
\begin{equation}
h \sim \varepsilon^{1/3}\,\max(1,|x_0|),
\end{equation}
where $\varepsilon$ is the relative machine precision. In multiprecision arithmetic, $\varepsilon$ depends on the working precision, so $h$ should adapt accordingly.

\subsection{Second-order (Taylor) propagation: mean correction and second-order variance term}

When nonlinearity of $f$ cannot be ignored and inputs can be approximated as \textbf{independent normal} random variables
$X_i\sim\mathcal{N}(\mu_i,\sigma_i^2)$, a second-order expansion around $\bm{\mu}$ gives
\begin{equation}
f(\bm{X})\approx f(\bm{\mu})+\bm{g}^\mathsf{T}\Delta\bm{x}
\;+\;\frac{1}{2}\Delta\bm{x}^\mathsf{T}H\,\Delta\bm{x},
\qquad \Delta\bm{x}=\bm{X}-\bm{\mu}.
\end{equation}
Here $\bm{g}=\nabla f$ and $H$ is the Hessian ($H_{ij}=\partial^2 f/\partial x_i\partial x_j$).

Under this approximation, the \textbf{mean} and \textbf{variance} can be written as
\begin{align}
\mathbb{E}[f(\bm{X})] &\approx f(\bm{\mu})+\frac{1}{2}\sum_{i=1}^{m} H_{ii}\,\sigma_i^2,\\
\mathrm{Var}[f(\bm{X})] &\approx \bm{g}^\mathsf{T}\bm{\Sigma}\bm{g}
\;+\;\frac{1}{2}\,\mathrm{tr}\!\left((H\bm{\Sigma})^2\right).
\end{align}
When $\bm{\Sigma}=\mathrm{diag}(\sigma_1^2,\ldots,\sigma_m^2)$ (independent inputs),
\begin{equation}
\mathrm{tr}\!\left((H\bm{\Sigma})^2\right)=\sum_{i=1}^{m}\sum_{j=1}^{m} H_{ij}^2\,\sigma_i^2\sigma_j^2.
\end{equation}

Notes:
\begin{itemize}
  \item Compared to first-order linearization, the second-order method can both correct mean shifts due to nonlinearity and include second-order Hessian contributions to the variance.
  \item The formulas rely on assumptions such as normality, small uncertainties, and smoothness. For non-Gaussian inputs or domain/discontinuity issues, Monte Carlo is recommended.
\end{itemize}

\subsection{Monte Carlo propagation}

Monte Carlo methods do not explicitly use derivatives; instead they sample from the input distributions:
\begin{equation}
X_i^{(k)}\sim \mathcal{N}(\mu_i,\sigma_i^2),\quad
Y^{(k)} = f\!\left(X_1^{(k)},\ldots,X_m^{(k)}\right),\quad k=1,\ldots,N.
\end{equation}
The output expectation and uncertainty are estimated by the sample mean and sample standard deviation:
\begin{equation}
\bar{Y}=\frac{1}{N}\sum_{k=1}^{N}Y^{(k)},\qquad
s_Y=\sqrt{\frac{1}{N-1}\sum_{k=1}^{N}\left(Y^{(k)}-\bar{Y}\right)^2}.
\end{equation}

Notes:
\begin{itemize}
  \item Naturally captures higher-order effects, nonlinearity, and distribution deformation due to truncation/domain constraints.
  \item Cost scales with sample count $N$; in practice fixed random seeds can be used for reproducibility.
  \item If samples hit invalid domains (e.g., $\sqrt{x}$ with $x<0$), the corresponding samples fail and must be discarded or modeled with a more appropriate input distribution.
\end{itemize}

\subsection{Root-solving uncertainty propagation}

For implicit root solving, DataLab supports independent-input first-order propagation through the implicit function theorem. Monte Carlo propagation samples uncertain constants and data inputs, solves each sampled problem, and reports the sample standard deviation of the resulting roots. Scalar second-order propagation is local and intended for smooth scalar real roots; for systems, branch-sensitive equations, or root-count changes, Monte Carlo is preferred.

\section{Fitting Methods (Fitting)}

\subsection{Least squares and weighted $\chi^2$}

Given data points $\{(x_i,y_i)\}_{i=1}^{n}$ and a model $y_{\mathrm{model}}(x;\bm{p})$, least squares fitting minimizes the residual sum of squares:
\begin{equation}
\mathrm{RSS}(\bm{p})=\sum_{i=1}^{n} r_i(\bm{p})^2,\qquad r_i=y_{\mathrm{model}}(x_i;\bm{p})-y_i.
\end{equation}
When each point has a standard uncertainty $\sigma_i$ and independent Gaussian noise is assumed, one often uses a weighted $\chi^2$:
\begin{equation}
\chi^2(\bm{p})=\sum_{i=1}^{n} w_i\,r_i(\bm{p})^2,\qquad w_i=\frac{1}{\sigma_i^2}.
\end{equation}
With degrees of freedom $\nu=n-k$ ($k$ free parameters), the reduced chi-square $\chi^2_\nu=\chi^2/\nu$ is often used to assess consistency of the noise model (ideally $\chi^2_\nu\approx 1$).

\subsection{Linear models: analytic solution and covariance}

When the model is linear in parameters:
\begin{equation}
y \approx \sum_{j=1}^{k} p_j\,\varphi_j(x),
\end{equation}
let the design matrix be $X_{ij}=\varphi_j(x_i)$. The weighted least-squares normal-equation solution is
\begin{equation}
\hat{\bm{p}}=(X^\mathsf{T}WX)^{-1}X^\mathsf{T}W\bm{y},
\end{equation}
where $W=\mathrm{diag}(w_1,\ldots,w_n)$. Under the classical approximation, the parameter covariance is
\begin{equation}
\mathrm{Cov}(\hat{\bm{p}})\approx \sigma^2 (X^\mathsf{T}WX)^{-1},\qquad
\sigma^2\approx \frac{\chi^2}{\nu}.
\end{equation}
The standard uncertainty of parameter $p_j$ is $\sigma(p_j)=\sqrt{\mathrm{Cov}_{jj}}$, and off-diagonal elements reflect parameter correlations.

\subsection{Nonlinear models: local linearization and Jacobian}

For nonlinear models, residuals can be linearized near the optimum:
\begin{equation}
\bm{r}(\bm{p}+\Delta\bm{p})\approx \bm{r}(\bm{p}) + J\,\Delta\bm{p},\qquad
J_{ij}=\frac{\partial r_i}{\partial p_j}.
\end{equation}
This leads to a Gauss--Newton covariance estimate:
\begin{equation}
\mathrm{Cov}(\hat{\bm{p}})\approx \sigma^2 (J^\mathsf{T}WJ)^{-1}.
\end{equation}
This is equivalent (up to a sign) to using the Jacobian of the model output $\partial y/\partial p$, and forms the basis of many uncertainty estimates in nonlinear fitting.

\subsection{Model comparison metrics: AIC / BIC / $R^2$}

To balance goodness-of-fit and complexity across models, information criteria are commonly used. Using $\mathrm{RSS}$ (or unweighted residual sum of squares) as a simplified noise scale:
\begin{align}
\mathrm{AIC} &\approx 2k + n\ln(\mathrm{RSS}/n),\\
\mathrm{BIC} &\approx k\ln n + n\ln(\mathrm{RSS}/n).
\end{align}
The coefficient of determination $R^2$ is often used to measure explained variance:
\begin{equation}
R^2 = 1-\frac{\sum_i r_i^2}{\sum_i (y_i-\bar{y})^2},
\end{equation}
but should be interpreted with care for nonlinear, no-intercept, or weighted cases.

\section{Statistical Averaging Methods (Statistics)}

\subsection{Arithmetic mean and standard error}

Arithmetic mean:
\begin{equation}
\bar{x}=\frac{1}{n}\sum_{i=1}^{n}x_i.
\end{equation}
Sample variance (unbiased estimate):
\begin{equation}
s^2=\frac{1}{n-1}\sum_{i=1}^{n}(x_i-\bar{x})^2,\qquad (n>1)
\end{equation}
Standard error of the mean:
\begin{equation}
\sigma_{\bar{x}}=\frac{s}{\sqrt{n}}.
\end{equation}

\subsection{Weighted mean (weights $w_i=1/\sigma_i^2$)}

When each point has a known standard uncertainty $\sigma_i$ (independent), the weighted mean is
\begin{equation}
\bar{x}_w=\frac{\sum_i w_i x_i}{\sum_i w_i},\qquad w_i=\frac{1}{\sigma_i^2}.
\end{equation}
Its standard uncertainty is
\begin{equation}
\sigma_{\bar{x}_w}=\sqrt{\frac{1}{\sum_i w_i}}.
\end{equation}
A commonly used effective sample size measures the concentration of weights:
\begin{equation}
n_{\mathrm{eff}}=\frac{\left(\sum_i w_i\right)^2}{\sum_i w_i^2},
\end{equation}
which may be much smaller than $n$ when weights vary strongly.

If one wants to estimate the scatter of the data under the same weights, a common unbiased weighted sample variance is
\begin{equation}
s_w^2=\frac{\sum_i w_i(x_i-\bar{x}_w)^2}{\sum_i w_i-\frac{\sum_i w_i^2}{\sum_i w_i}},
\end{equation}
which reduces to the ordinary sample variance $s^2$ when all weights are equal. The denominator can also be written as $W(1-1/n_{\mathrm{eff}})$ where $W=\sum_i w_i$.

\section{Practical Tips (method-agnostic)}
\begin{itemize}
  \item Extrapolation/acceleration results should be cross-validated against the original sequence trend, differences, and multiple methods.
  \item If significant correlations exist in uncertainty propagation, provide a covariance matrix or use Monte Carlo.
  \item Fitting uncertainties depend on model correctness and noise assumptions; for strongly correlated parameters, inspect covariance/correlation coefficients.
  \item Weighted statistics require credible $\sigma_i$; avoid treating systematic errors as purely statistical uncertainties.
\end{itemize}

\end{document}
